\section{Volitional Energy}

\begin{definition}[Volitional energy]
Volitional energy $\Will(t)$ is the limited capacity available at time $t$ to
interrupt the current behavioural frame and initiate stabilizing action toward
a target frame.
\end{definition}

\begin{discussion}
The practical value of this definition is not that it gives the reader a
precise biological meter of ``will.'' Its value is architectural. The agent
has only a limited ability to force a transition by direct effort, while the
barrier separating the current frame from the target frame can fluctuate
substantially across the day. This asymmetry matters because most people can
improve timing more reliably than they can summon heroic motivation on demand.
\end{discussion}

\begin{definition}[Leverage]
The effective leverage of free will at time $t$ is
\[
  \Leverage(t) = \frac{\Will(t)}{\DV(t)}.
\]
\end{definition}

\begin{proposition}[Timing dominates intensity under bounded will]
If short-horizon increases in $\Will(t)$ are small relative to ordinary
fluctuations in $\DV(t)$, then the most reliable way to raise
$\Leverage(t)$ is to act when the barrier is already falling rather than to
wait for a dramatic increase in motivation.
\end{proposition}

\begin{proof}[Proof sketch]
By definition, leverage improves when the numerator rises or the denominator
falls. In ordinary life, abrupt large gains in available will are rare and
costly to sustain, whereas barrier changes are common at wake-up, task
completion, arrival in a new location, or after a strong informational update.
Therefore a timing strategy that exploits barrier troughs usually dominates a
strategy that depends on exceptional effort.
\end{proof}

\begin{theorem}[Decision-window principle]
Free will has maximum strategic effect not when effort is most intense, but
when the activation barrier is locally minimal relative to available
volitional energy.
\end{theorem}

\begin{proof}[Proof sketch]
If leverage is represented by $\Will(t)/\DV(t)$, then leverage rises when
$\Will(t)$ increases or $\DV(t)$ decreases. In ordinary life, large increases
in available will are unreliable, but temporary drops in barrier are common:
after waking, after finishing a task, upon entering a new environment, or
immediately after a sharp emotional update. Therefore the most reliable
strategy is to act during barrier troughs rather than waiting for heroic
motivation.
\end{proof}

\begin{example}[Why raw sincerity often fails]
A student may genuinely value mathematics at night and still fail to begin a
proof-heavy session if the current frame is stabilized by fatigue, clutter,
and open-ended task ambiguity. The same student may start successfully the
next morning with less subjective passion because the barrier has fallen: the
old frame is weaker, the first task is already visible, and the cost of
redirecting attention is lower.
\end{example}
